% ClearVision: Real-Time DIBR for A-Pillar Blind Spot Elimination
% Seán Kelly — 21421506 — University of Galway — 1MECE1

\documentclass[11pt, a4paper]{report}

% -----------------------------------------------------------------------
%  Packages
% -----------------------------------------------------------------------
\usepackage[a4paper, top=2.5cm, bottom=2.5cm, left=2.8cm, right=2.8cm]{geometry}
\usepackage{graphicx}
\usepackage{amsmath}
\usepackage{booktabs}
\usepackage{hyperref}
\usepackage{listings}
\usepackage{xcolor}
\usepackage{caption}
\usepackage{parskip}
\usepackage{fancyhdr}
\usepackage{titlesec}
\usepackage{enumitem}
\usepackage{microtype}

% -----------------------------------------------------------------------
%  Placeholder figure command
% -----------------------------------------------------------------------
\newcommand{\figplaceholder}[2][0.82\linewidth]{%
  \begin{center}
  \fbox{\begin{minipage}{#1}\centering\vspace{2.2cm}
  \textit{[Figure placeholder: #2]}\vspace{2.2cm}
  \end{minipage}}
  \end{center}
}

% -----------------------------------------------------------------------
%  Listings style (GLSL / C++ / Python snippets)
% -----------------------------------------------------------------------
\lstset{
    basicstyle=\ttfamily\small,
    breaklines=true,
    frame=single,
    captionpos=b,
    commentstyle=\color{gray},
    keywordstyle=\color{blue!70!black},
    stringstyle=\color{red!60!black},
    showstringspaces=false,
    tabsize=4,
    numbers=left,
    numberstyle=\tiny\color{gray},
    numbersep=5pt,
}

% -----------------------------------------------------------------------
%  Header / Footer
% -----------------------------------------------------------------------
\pagestyle{fancy}
\fancyhf{}
\fancyhead[L]{\small ClearVision --- Project Report}
\fancyhead[R]{\small Se\'{a}n Kelly, 21421506, 1MECE1}
\fancyfoot[C]{\thepage}
\renewcommand{\headrulewidth}{0.4pt}

% -----------------------------------------------------------------------
%  Title formatting
% -----------------------------------------------------------------------
\titleformat{\chapter}[block]
  {\normalfont\Large\bfseries}{\thechapter.}{0.8em}{}
\titlespacing*{\chapter}{0pt}{20pt}{10pt}

\setlength{\parindent}{0pt}
\setlength{\parskip}{6pt}

% -----------------------------------------------------------------------
%  Document
% -----------------------------------------------------------------------
\begin{document}

% =====================================================================
%  COVER PAGE
% =====================================================================
\begin{titlepage}
    \centering
    \vspace*{1.5cm}

    {\Large \textsc{University of Galway}}\\[0.4em]
    {\large School of Engineering}\\[0.3em]
    {\large Department of Electronic and Electrical Engineering}\\[2.8cm]

    \rule{\linewidth}{1pt}\\[0.5em]
    {\LARGE \bfseries ClearVision: A Real-Time Vision\\[0.3em]
    System for the A-Pillar Blind Spot}\\[0.5em]
    \rule{\linewidth}{1pt}\\[2cm]

    {\large \textbf{Author:} Se\'{a}n Kelly}\\[0.4em]
    {\large \textbf{Student ID:} 21421506}\\[0.4em]
    {\large \textbf{Date:} 11 May 2026}\\[0.4em]
    {\large \textbf{Supervisor:} Prof.\ Martin Glavin}\\[0.4em]
    {\large \textbf{Co-Assessor:} Dr.\ Brian Deegan}\\[0.4em]

    \newpage

    \begin{minipage}{0.85\linewidth}
    \begin{abstract}
    \noindent
    The structural A-pillars at the corners of a vehicle's windscreen create
    a persistent blind spot that obstructs the driver's view of the road ahead.
    This report describes ClearVision, a system designed to make this obstruction
    effectively transparent: a display panel mounted in the pillar shows a
    synthesised view of the hidden scene, updated in real time to track the
    driver's head position, so that the panel behaves like a window rather than
    a screen.

    The system captures the scene using a Luxonis OAK-D Lite stereo camera,
    which provides hardware-accelerated depth estimation.
    A viewer-facing USB webcam tracks the driver's head position using a
    lightweight face-mesh pipeline.
    A multi-pass compute shader pipeline on an NVIDIA Jetson Nano then
    reprojects the depth-coloured scene to match the driver's instantaneous
    viewpoint, producing physically-correct motion parallax.

    The project evolved through three development phases, achieving
    25--30~FPS at 640$\times$480 in the final system, compared to
    approximately 5~FPS in the initial prototype.
    Key contributions include a geometrically-correct virtual window
    projection model and a Jump Flood Algorithm hole-fill pass for
    disocclusion regions.
    \end{abstract}
    \end{minipage}
\end{titlepage}

% =====================================================================
%  TABLE OF CONTENTS
% =====================================================================
\tableofcontents
\listoftables
\listoffigures
\newpage

% =====================================================================
%  CHAPTER 1 — INTRODUCTION
% =====================================================================
\chapter{Introduction}

\section{Motivation}

The structural A-pillars at the front corners of a vehicle obstruct
approximately 4--8$^\circ$ of the driver's forward field of
view~\cite{nhtsa_visibility}.
Although this angle appears small, at a pedestrian crossing 20~m ahead it
translates to roughly 1.4~m of hidden lateral space --- sufficient to
conceal a cyclist or a child.
The problem is particularly acute at low-speed urban junctions, where
pedestrians and cyclists frequently approach from the front quarter.

\begin{figure}[h]
\centering
\figplaceholder{A-pillar blind spot concept. Left: the driver's
  field of view, showing the angular region obscured by the A-pillar.
  Right: the ClearVision concept --- a display mounted in the pillar
  renders a viewpoint-tracked synthesised image, so that the pillar
  appears transparent to the driver.}
\caption{The A-pillar blind spot and the ClearVision concept.}
\label{fig:apillar_concept}
\end{figure}

Common mitigation strategies include convex mirrors and fixed-viewpoint
camera-display systems.
Convex mirrors increase field of view at the cost of significant spatial
distortion, making accurate distance judgement difficult.
Fixed camera displays, fitted to an increasing number of production
vehicles, improve coverage but suffer from a fundamental perceptual
limitation: because the displayed image does not shift as the driver
moves, the brain cannot use motion parallax to resolve depth.
Objects shown on a static camera feed appear flat, making it hard to judge
their distance or speed.

ClearVision takes a different approach.
Rather than displaying a static camera feed, it synthesises a novel
viewpoint that tracks the driver's head position.
The display panel appears transparent: objects in the hidden scene move
naturally as the driver shifts position, exactly as they would through a
real window.
This preserves the depth cues that drivers rely on for accurate object
localisation.

\section{Prior Approaches to the A-Pillar Problem}

Several approaches to A-pillar transparency have been explored in industry
and research.

\textbf{Camera-display systems} are the most widely deployed approach.
A camera mounted on the exterior of the pillar feeds a display embedded in
the interior face.
This is mechanically straightforward and requires no depth processing.
Its key limitation is that the displayed image is a fixed viewpoint:
the driver's head position has no effect on what is shown.
This breaks the motion parallax cues that the brain uses to judge depth
and speed, and requires the driver to consciously interpret a 2D image
rather than directly perceiving the scene.

\textbf{Transparent pillar concepts} attempt to make the physical pillar
optically transparent.
Toyota demonstrated a prototype in 2014 using OLED panels affixed to the
exterior surface of the pillar, fed by a camera behind it~\cite{toyota_apillar}.
Continental presented a similar concept using a curved OLED display
integrated into the pillar cross-section~\cite{continental_apillar}.
Jaguar Land Rover's 360 Virtual Urban Windscreen concept (2014) embedded
screens in the A-, B-, and C-pillars fed by exterior cameras, with the
display activating automatically when the system detected a lane change
or junction approach~\cite{jlr2014pillar}.
All of these systems share the fixed-viewpoint limitation: because the
displayed image is a live camera feed from a fixed position, head movement
produces no parallax, and the driver cannot use depth perception to
judge the distance or speed of objects in the hidden region.

\textbf{Viewpoint-adaptive systems} address this by updating the displayed
image in response to the driver's position.
This is the approach taken by ClearVision.
The technical challenge is generating a geometrically correct novel
viewpoint in real time from standard camera input, which requires a depth
map of the scene and a rendering pipeline fast enough to run at video rate.

\section{Related Techniques}
\label{sec:related}

\subsection{View Synthesis}

The core technical problem --- generating a novel camera viewpoint from
existing images --- has been studied extensively in the computer vision
and graphics literature.

Depth-Image Based Rendering (DIBR), introduced for 3DTV
applications~\cite{fehn2004depth,smolic2007coding}, synthesises new
viewpoints by warping each pixel of a reference image to the position it
would occupy in the target view, guided by a depth map.
It operates in real time without requiring a full 3D model of the scene,
making it suitable for embedded deployment.
The key limitation is disocclusion: regions visible from the target
viewpoint but not from the reference camera produce holes that must be
filled heuristically.
M\"{u}ller et al.~\cite{muller2008viewsynth} survey disocclusion handling
methods for 3DTV view synthesis, including background depth propagation
and inpainting; the Jump Flood Algorithm~\cite{rong2006jfa} offers an
efficient GPU-parallel alternative for 2D Voronoi hole fill.

More recent approaches, such as Neural Radiance Fields
(NeRF)~\cite{mildenhall2020nerf}, produce substantially higher quality
novel views but require minutes to hours of per-scene training and
significant GPU compute for inference, making them unsuitable for
real-time embedded use.
DIBR remains the practical choice for this application.

\subsection{Stereo Depth Estimation}

A reliable depth map is central to DIBR.
Stereo vision --- computing depth from the disparity between two
horizontally offset cameras --- is the dominant approach for real-time
embedded use~\cite{scharstein2002taxonomy}.
Block Matching (BM) and Semi-Global Block Matching
(SGBM)~\cite{hirschmuller2005sgbm} are the primary classical algorithms.
SGBM produces significantly better depth maps than BM, particularly at
object boundaries and in textured regions, at the cost of higher
computation.
On a CPU, SGBM at 640$\times$480 typically requires 20--30~ms --- too
slow to leave headroom for the rendering pipeline on a Jetson Nano.

Hardware-accelerated stereo, as provided by the Luxonis OAK-D
Lite~\cite{luxonis_oakd}, offloads this computation to dedicated silicon
on the Myriad X VPU, returning subpixel disparity at the camera frame
rate with no host CPU cost.
Structured-light and time-of-flight sensors are alternatives that perform
well on textureless surfaces, but at substantially higher cost and with
reduced operating range in direct sunlight.

\subsection{Head and Face Tracking}

Estimating the driver's 3D head position in real time requires a face
detection and tracking pipeline that can run on the remaining compute budget.
Classical approaches --- Haar cascade detectors~\cite{viola2004robust}
and optical flow trackers --- are lightweight but provide only 2D screen
coordinates, requiring an additional depth cue.

Deep learning face mesh models, such as the MediaPipe
FaceMesh~\cite{kartynnik2019facemesh}, provide dense 3D facial landmark
estimates from a monocular camera.
The iris landmarks are particularly useful: the human iris has a
near-constant physical diameter of approximately 11.7~mm across the adult
population~\cite{bergmanson2017iris}, allowing monocular depth estimation
directly from its projected size, without a separate depth sensor on the
driver-facing camera.
TensorRT acceleration makes this feasible on the Jetson Nano's Maxwell GPU
at the required frame rate.

\subsection{Head-Coupled Displays and Off-Axis Projection}

The idea of updating a rendered viewpoint to match the viewer's physical
position relative to a flat screen --- creating a ``window'' into a
virtual world --- has a long research history.
Ware, Arthur, and Booth~\cite{ware1993fishtank} formalised this as
\emph{fish tank virtual reality} (FTVR), demonstrating in 1993 that
head-coupled stereo reduced 3D object-placement error from 22\% to 1.3\%
compared to a fixed display.
Their key finding was that head coupling contributed more to depth
perception than stereoscopy alone.

For scenes with known 3D geometry --- synthetic environments in game
engines or interactive visualisations --- view synthesis reduces to
repositioning the virtual camera: no depth estimation, no hole fill.
Lee's Wii Remote demonstration~\cite{lee2008head} made this accessible
to a general audience, coupling a head-tracked IR sensor to an OpenGL
scene on a standard monitor.
Kooima~\cite{kooima2009generalized} formalised the off-axis asymmetric
frustum construction that underpins these systems.

ClearVision faces a harder version of the same problem: the scene is a
real-world camera feed, not a 3D model.
A depth map must be estimated, the scene reprojected to the novel
viewpoint, and disocclusion holes filled --- steps that do not exist
when the geometry is already known.
This is the distinction between interactive graphics and image-based
rendering, and is the central technical challenge of the project.

\section{Project Scope and Objectives}

The goal of ClearVision is to demonstrate that a physically-correct,
real-time A-pillar transparency system is achievable on commodity embedded
hardware, using passive stereo depth and monocular head tracking.

The system is designed to:

\begin{enumerate}[noitemsep]
    \item Capture the scene outside the vehicle using a stereo camera
          mounted at the A-pillar position.
    \item Estimate the driver's 3D head position in real time from a
          viewer-facing monocular camera.
    \item Synthesise a novel viewpoint of the captured scene corresponding
          to the driver's current position, so that the display behaves as
          a geometrically correct transparent window.
    \item Achieve a display frame rate sufficient for smooth, perceptible
          parallax motion (target: $\geq$20~FPS at 640$\times$480).
\end{enumerate}

The system targets the NVIDIA Jetson Nano as the host compute platform,
representing the class of low-power embedded GPU hardware available in
automotive contexts.
The evaluation focuses on frame rate, depth quality, and the perceptual
quality of the parallax effect.
Full vehicle integration, multi-driver calibration, and weatherproofing
are out of scope for this project.

% =====================================================================
%  CHAPTER 2 — TECHNICAL BACKGROUND
% =====================================================================
\chapter{Technical Background}

This chapter introduces the key technical concepts underpinning the
ClearVision system.
The rendering pipeline and implementation details are covered in
Chapters~5--7; this chapter provides the theoretical foundation.

\section{Depth-Image Based Rendering}
\label{sec:dibr}

Depth-Image Based Rendering (DIBR) synthesises novel camera viewpoints
from a reference image and its associated per-pixel depth map, without
requiring a full 3D scene model~\cite{fehn2004depth}.
Each pixel in the reference image is unprojected to a 3D world point
using the depth value and camera intrinsics, then reprojected into the
target viewpoint to determine its position in the output image.

A stereo camera pair provides the depth map.
The two cameras are separated by a known \textbf{baseline} $B$ (metres).
For a scene point at depth $Z$, the horizontal pixel offset between the
two images --- the \textbf{disparity} $d$ --- satisfies:

\begin{equation}
    d = \frac{f \cdot B}{Z}
    \label{eq:disparity}
\end{equation}

where $f$ is the focal length in pixels.
Inverting, given a measured disparity $d$:

\begin{equation}
    Z = \frac{f \cdot B}{d}
    \label{eq:depth_from_disparity}
\end{equation}

Given $Z$ and camera intrinsics $(f_x, f_y, c_x, c_y)$, a pixel $(u, v)$
is unprojected to metric 3D coordinates:

\begin{equation}
    X = \frac{(u - c_x) \cdot Z}{f_x}, \quad
    Y = \frac{(v - c_y) \cdot Z}{f_y}
    \label{eq:unproject}
\end{equation}

These world-space points form a point cloud.
Any novel viewpoint can be rendered by projecting the point cloud into
the target camera frame.

\section{The Virtual Window Model}
\label{sec:vwindow}

A naive DIBR implementation treats view synthesis as a horizontal pixel
shift: $u' = u - \alpha \cdot d$, where $\alpha$ blends between the two
stereo camera positions.
For the A-pillar application, this has two fundamental problems.
First, the viewer is \emph{inside} the car while the cameras are
\emph{outside}; the viewer is not interpolating between the stereo pair,
so no value of $\alpha$ corresponds to a physically meaningful viewpoint.
Second, the shift has no connection to physical viewer displacement ---
the scale constant must be hand-tuned empirically.

The virtual window model replaces this with a geometrically correct ray
projection.
For each output pixel on the display, a ray is cast from the driver's
head position through the display surface into the scene.
The ray is intersected with the display plane to find the pixel's position
on the display, then followed into the point cloud to determine its colour.

Concretely, for each world point $W$ and head position $H$:

\begin{enumerate}[noitemsep]
    \item Compute the ray direction:
          $\mathbf{r} = \mathrm{normalize}(W - H)$.
    \item Intersect $\mathbf{r}$ with the display plane (defined by a
          point $P$ and outward normal $\hat{n}$):
          \begin{equation}
              t = \frac{(P - H) \cdot \hat{n}}{r \cdot \hat{n}}, \quad
              Q = H + t \cdot r
              \label{eq:ray_plane}
          \end{equation}
    \item Express $Q$ in display-local UV coordinates to determine which
          output pixel receives this world point's colour.
\end{enumerate}

This yields physically correct parallax by construction, with no tuning
constants.
Horizontal and vertical parallax both emerge naturally from the geometry.

\section{Off-Axis Projection}
\label{sec:offaxis}

The virtual window model is closely related to off-axis (asymmetric
frustum) projection~\cite{kooima2009generalized}, demonstrated
interactively by Lee~\cite{lee2008head}.
In this formulation, the OpenGL viewing frustum is defined to match the
viewer's physical position relative to the screen, so that the rendered
synthetic scene exhibits correct parallax as the viewer moves.

ClearVision incorporates this in two forms.
Early proof-of-concept demonstrations used an orthographic asymmetric
frustum in Python to validate the parallax effect.
The final system generalises this to full 3D with real depth data,
implemented as a compute shader pipeline described in Chapter~7.

\section{Stereo Vision and Disparity Estimation}
\label{sec:stereo}

Reliable disparity estimation requires \textbf{stereo rectification}:
both camera images are re-projected onto a common plane so that
corresponding scene points lie on the same image row.
After rectification, a matching algorithm searches along horizontal
scanlines for the best-matching pixel offset.

Semi-Global Block Matching (SGBM)~\cite{hirschmuller2005sgbm} is the
standard algorithm for high-quality stereo matching.
It aggregates matching costs along multiple scanline directions, producing
significantly better results than simple block matching at object
boundaries.
Classical stereo algorithms are surveyed by Scharstein and
Szeliski~\cite{scharstein2002taxonomy}.

The OAK-D Lite performs rectification and SGBM in hardware on its Myriad
X VPU, returning subpixel-accurate disparity at the camera frame rate
with zero host CPU cost.

\section{Face Tracking and Iris-Based Depth}
\label{sec:facetrack}

The system uses a two-stage face tracking pipeline to estimate the
driver's 3D head position.
A Haar cascade detector~\cite{viola2004robust} provides a face bounding
box; a deep learning face mesh model~\cite{kartynnik2019facemesh} then
refines this to 478 3D facial landmarks, including iris landmarks.

Depth is recovered from the apparent iris diameter.
The human iris has a near-constant physical diameter of approximately
11.7~mm~\cite{bergmanson2017iris} (population variance $\pm 0.5$~mm ---
substantially tighter than inter-ocular distance or face width).
Its projected size in pixels, $d_{\mathrm{iris,px}}$, depends on focal
length $f_t$ and distance $Z$ via the pinhole model:

\begin{equation}
    Z = \frac{D_{\mathrm{iris}} \cdot f_t}{d_{\mathrm{iris,px}}}
    \label{eq:iris_depth}
\end{equation}

Lateral and vertical head position follow directly from the normalised
landmark coordinates once $Z$ is known.

% =====================================================================
%  CHAPTER 3 — SYSTEM ARCHITECTURE
% =====================================================================
\chapter{System Architecture}

\section{Hardware}
\label{sec:hardware}

Three hardware components comprise the system.

\textbf{OAK-D Lite (Luxonis)}~\cite{luxonis_oakd} is the external stereo
camera, mounted at the A-pillar position facing outward.
It contains two OmniVision OV7251 global-shutter monochrome sensors
($640\times480$, 73$^\circ$ HFOV) on a 75~mm fixed baseline, a Sony
IMX214 colour camera, and an Intel Myriad X VPU.
The global shutter is important: a rolling shutter introduces lateral skew
artefacts during camera or scene motion, distorting the disparity map.
The OAK-D Lite connects via a single USB-C cable and is bus-powered.

\textbf{NVIDIA Jetson Nano (B01)} is the host compute platform,
providing a Maxwell GPU (128 CUDA cores, OpenGL ES~3.1) running
NVIDIA JetPack~4.x (L4T~R32).
The Maxwell GPU is sufficient for the DIBR compute shader pipeline
once the SGBM bottleneck is offloaded to the OAK-D VPU.

\textbf{USB webcam} (interior, viewer-facing) feeds the head tracking
pipeline.
No stereo hardware is required on this side; a single camera with the
iris depth estimator suffices.

\begin{figure}[h]
\centering
\figplaceholder{Hardware assembly diagram. OAK-D Lite mounted externally
  at the A-pillar position (inverted, facing outward); USB-C connection
  to the Jetson Nano host; viewer-facing USB webcam for head tracking.}
\caption{ClearVision hardware assembly.}
\label{fig:hardware_rig}
\end{figure}

\section{Architecture Overview}
\label{sec:arch_overview}

The system was developed in three phases, each addressing a distinct
set of limitations.
Table~\ref{tab:arch_evolution} summarises the key changes.

\begin{table}[h]
\centering
\caption{Architecture evolution across development phases}
\label{tab:arch_evolution}
\begin{tabular}{@{}p{3.5cm}p{3.5cm}p{3.5cm}@{}}
\toprule
\textbf{Phase 1} & \textbf{Phase 2} & \textbf{Phase 3 (Final)} \\
\midrule
Dual CSI cameras & Dual CSI cameras & OAK-D Lite (USB-C) \\
CPU SGBM / CUDA BM & CPU SGBM / CUDA BM & Hardware SGBM on Myriad X \\
Manual stereo calibration & Manual stereo calibration & EEPROM intrinsics \\
Disparity shift model & Disparity shift model & Virtual window ray model \\
Optical flow / OpenCV DNN & MOSSE + Haar Cascade & TRT FaceMesh + iris depth \\
Horizontal hole fill & Horizontal hole fill & JFA 2D hole fill \\
$\sim$5 FPS & $\sim$10 FPS & 25--30 FPS \\
\bottomrule
\end{tabular}
\end{table}

Phase~1 established the core DIBR pipeline using dual CSI cameras
connected directly to the Jetson Nano.
Stereo rectification and matching ran on the host CPU, and SGBM dominated
the frame budget at 20--30~ms per frame.
Phase~2 introduced head tracking, using a MOSSE correlation filter
anchored by a Haar cascade detector.
This achieved acceptable CPU load, but the tracker could not provide
metric depth, which is required to correctly scale the parallax in the
virtual window model.
Phase~3 replaced the CSI rig with an OAK-D Lite, offloading stereo
entirely to the Myriad X VPU, and migrated both the projection model
and the head tracker to metric-aware implementations.

\section{Software Architecture}
\label{sec:software_arch}

The system is implemented in C\texttt{++} using Qt for the OpenGL window
and event loop.
Three concurrent threads handle the three data sources:

\begin{itemize}[noitemsep]
    \item \textbf{OAK receiver thread} (\texttt{oak\_receiver.cpp}):
          polls the DepthAI device for disparity (RAW16 subpixel),
          rectified mono frames, and colour frame; writes to shared
          \texttt{cv::Mat} objects protected by a mutex.
    \item \textbf{Face tracker thread} (\texttt{face\_tracker.cpp}):
          runs Haar cascade detection and TensorRT FaceMesh inference on
          the webcam feed; outputs a smoothed 3D head position vector.
    \item \textbf{Qt render thread}: calls \texttt{paintGL()} at the
          OAK-D frame rate, consuming the latest data from both background
          threads and dispatching the compute shader pipeline.
\end{itemize}

\begin{figure}[h]
\centering
\figplaceholder{Software architecture diagram. Three threads feed the
  Qt render thread: OAK receiver (disparity, rectified frames, colour),
  face tracker (head position vector), and the render thread itself
  (compute shader dispatch, display output). Shared data protected by
  mutexes.}
\caption{ClearVision software architecture.}
\label{fig:software_architecture}
\end{figure}

\section{Compute Shader Pipeline}
\label{sec:shader_pipeline}

Each frame, the following GPU passes execute in sequence:

\begin{table}[h]
\centering
\caption{Compute shader passes per frame}
\label{tab:shader_passes}
\begin{tabular}{@{}lll@{}}
\toprule
Pass & Shader & Function \\
\midrule
0 & \texttt{CLEAR\_CS}          & Zero the depth SSBO (GPU-side) \\
1 & \texttt{UNPROJECT\_CS}      & Disparity $\to$ world-space XYZ texture \\
2 & \texttt{VWINDOW\_DEPTH\_CS}  & World points $\to$ virtual window depth buffer \\
3 & \texttt{VWINDOW\_COLOR\_CS}  & Depth-tested colour splat from left camera \\
4 & \texttt{VWINDOW\_RIGHT\_CS}  & Right-camera disocclusion fill \\
5 & \texttt{JFA\_INIT\_CS}      & Seed JFA with filled pixel coordinates \\
6--N & \texttt{JFA\_CS} $\times \lceil\log_2(\max(W,H))\rceil$ & Voronoi propagation \\
N+1 & \texttt{JFA\_GATHER\_CS}  & Copy nearest-seed colour into holes \\
Display & \texttt{DISPLAY\_VS/FS} & Fullscreen quad, letterboxed \\
\bottomrule
\end{tabular}
\end{table}

All image-space compute shaders use a 16$\times$16 workgroup layout.

\section{Stereo Baseline Selection}
\label{sec:baseline}

The OAK-D Lite's fixed 75~mm baseline is well-matched to the
application.
The A-pillar on a typical passenger vehicle occludes approximately
80--120~mm of lateral space at the driver's eye position.
For the virtual window parallax to correctly reveal the hidden scene,
the synthesisable viewpoint shift must span this same distance.

At the target depth range of 3--5~m, the OAK-D stereo pair
($f \approx 381$~px at 480P, 73$^\circ$ HFOV) produces disparities in
the range 5.7--9.5~px (computed from Equation~\ref{eq:disparity}),
both above the hardware SGBM noise floor.
Benchmarks from~\cite{arxiv2501} report $<$2\% depth error at 4~m for
OAK-D hardware under controlled conditions.
A detailed derivation and discussion of the baseline constraints is given
in Appendix~\ref{app:baseline}.

% =====================================================================
%  CHAPTER 4 — PROOF-OF-CONCEPT DEMONSTRATIONS
% =====================================================================
\chapter{Proof-of-Concept Demonstrations}

Before committing to a full C\texttt{++} implementation, two Python
proof-of-concept demonstrations were developed to validate the core
geometry and identify design issues early.
Python was chosen for its rapid iteration cycle: no compilation step,
direct OpenGL access via PyOpenGL, and a mature MediaPipe Python API.
The two demonstrations together established the key design decisions ---
the threading pattern, the sign conventions, and the projection
formulation --- that the C\texttt{++} pipeline subsequently inherited.

\section{Head Pose Vector Demo}

The first demonstration (\texttt{face\_landmarker\_demo.py}) validated
the face tracking foundation.
The goal was to render the 3D face orientation vector overlaid on a live
camera feed, without relying on any hardcoded scene constants.
This established that the MediaPipe FaceLandmarker API could be operated
in real time on a laptop, and that the landmark geometry was stable and
usable as a tracking signal.

The face normal vector is computed from a small set of facial landmarks ---
nose tip, chin, and inner eye corners --- using a cross product:

\begin{enumerate}[noitemsep]
    \item Eye line vector: $\mathbf{e} = P_{\mathrm{right}} - P_{\mathrm{left}}$
    \item Vertical face axis: $\mathbf{v} = P_{\mathrm{nose}} - P_{\mathrm{chin}}$
    \item Face normal: $\hat{n} = \mathrm{normalize}(\mathbf{v} \times \mathbf{e})$
\end{enumerate}

MediaPipe's LIVE\_STREAM mode was used throughout.
This mode delivers results asynchronously via a callback, which must not
perform rendering (OpenGL contexts are not thread-safe).
The callback stores results into a shared global protected by a lock;
the main thread reads the results and renders on each frame.
This producer-consumer pattern became the template for all subsequent
pipelines, including the final C\texttt{++} face tracker and OAK-D
receiver threads.

\begin{figure}[h]
\centering
\figplaceholder{Head pose vector demo screenshot. The 3D face normal
  vector is rendered as an arrow overlaid on the live webcam feed.
  The vector rotates smoothly as the face moves, with no hardcoded
  scaling constants.}
\caption{Head pose vector proof-of-concept demonstration.}
\label{fig:demo_head_pose}
\end{figure}

\section{Off-Axis Projection Demo}

The second demonstration (\texttt{head\_coupled\_offaxis.py}) coupled
the viewer's 2D nose-tip position from MediaPipe to an OpenGL rendered
scene.
Its purpose was to verify that a physically meaningful parallax effect
was perceptually convincing before the full depth-based system was built.

A calibration phase records the reference nose-tip position at the
viewer's normal viewing distance.
During the demo, the displacement from this reference drives the OpenGL
projection:

\begin{equation}
    \Delta x_{\mathrm{mm}} = -(x_{\mathrm{nose}} - x_{\mathrm{ref}}) \cdot W_{\mathrm{screen,mm}}
    \label{eq:delta_x}
\end{equation}

A negative sign is required because the webcam image is mirrored: a
rightward head movement appears as a leftward nose shift in camera
coordinates.
An orthographic frustum shifted by the head offset produces the parallax
effect:

\begin{lstlisting}[language=Python,
  caption={Orthographic frustum with head-coupled offset.
    The frustum origin shifts by the head displacement, so the scene
    appears to rotate in the opposite direction --- the parallax effect.}]
zoom = 2.0          # field of view factor
scale = 2.0         # movement amplification

glOrtho(-W/2 * zoom - dx * scale,
         W/2 * zoom - dx * scale,
        -H/2 * zoom - dy * scale,
         H/2 * zoom - dy * scale,
        -1000, 1000)
\end{lstlisting}

The orthographic formulation was adopted after the perspective frustum
failed to produce a visible effect: the combination of a 1~mm near plane
and millimetre-scale scene geometry placed all scene content behind the
near clipping plane.
The orthographic version demonstrated the head-coupled parallax effect
convincingly, validating the approach before the C\texttt{++} pipeline
was built.

\begin{figure}[h]
\centering
\figplaceholder{Off-axis projection demo screenshots. Left: viewer
  at the reference position --- the rendered cube is centred on screen.
  Right: viewer shifted laterally --- the frustum shifts to expose the
  scene region behind the new sightline, producing visible motion
  parallax. The effect is perceptually convincing at displacements
  above approximately 20~mm.}
\caption{Off-axis projection proof-of-concept demonstration.}
\label{fig:demo_offaxis}
\end{figure}

\subsection{Development Challenges}

Several non-trivial issues were encountered and resolved during
development of the demos.

\textbf{Texture upload failure on Windows.}
The demo was developed on Windows and verified on a Raspberry Pi.
On Windows, \texttt{glTexImage2D} failed silently when passed a
JPEG-loaded NumPy array; switching the source image to PNG resolved the
issue.
The Raspberry Pi's OpenGL driver was unaffected.
This highlighted the importance of testing on the target platform early.

\textbf{Axis inversions.}
MediaPipe returns nose coordinates in a mirrored camera frame where $+Y$
is downward.
After calibration, the head offset required sign inversion on both axes:
$\Delta x \to -\Delta x$ (camera $+X$ is the viewer's left) and
$\Delta y \to +\Delta y$ (prior incorrect negation removed).
These sign conventions were carried forward directly into the final
C\texttt{++} coordinate mapping (Section~\ref{sec:coord_mapping}).

\textbf{LIVE\_STREAM threading.}
IMAGE mode processes one frame at a time and was 5--10$\times$ slower
than LIVE\_STREAM mode.
Discovering this early, and establishing the correct threading model for
MediaPipe callbacks, avoided a significant re-architecture later.

\section{Value of the Demonstrations}

The two demonstrations, while simple, provided substantial confidence in
the chosen approach before any C\texttt{++} code was written.

Demo~1 established that MediaPipe's face landmark API was reliable enough
for real-time use, that the face normal computation was stable, and that
the producer-consumer threading pattern worked correctly.
Demo~2 established that the parallax effect was perceptually effective at
head displacements typical of a seated driver ($\pm$30--80~mm), that the
orthographic projection was the correct formulation for the geometry, and
that all the sign conventions in the coordinate system were identified and
correct before the full 3D implementation was attempted.

Together, they substantially reduced the implementation risk of the
C\texttt{++} pipeline and meant that the full system could be built with
confidence in the underlying approach.

% =====================================================================
%  CHAPTER 5 — THE DIBR PIPELINE: STEREO DEPTH
% =====================================================================
\chapter{The DIBR Pipeline: Stereo Depth}

\section{OAK-D Lite Integration}

\subsection{Replacing the Initial Stereo Rig}

The Phase~1 stereo pipeline used two IMX219 CSI cameras connected
directly to the Jetson Nano, with rectification and SGBM running
entirely on the host.
CPU SGBM dominated the frame budget at 20--30~ms per frame; even with
CUDA StereoBM this fell to 10--15~ms, but consumed GPU resources needed
by the DIBR render pipeline.
A manual stereo calibration procedure (\texttt{calibration.cpp}) was
also required each time the rig was adjusted.

The OAK-D Lite replaces the dual-CSI rig with a single USB-C connection.
The Myriad X VPU performs rectification and SGBM entirely in hardware,
returning already-rectified frames and subpixel disparity at zero host
CPU cost.
Camera intrinsics and baseline are read from the device EEPROM via
\texttt{device->readCalibration()}, eliminating the manual calibration
step.

\subsection{DepthAI Pipeline Configuration}

The pipeline graph is configured at startup.
Four DepthAI nodes are used: left and right MonoCamera (480P), a
StereoDepth node, and a ColorCamera (ISP output).
The stereo node is configured with:

\begin{itemize}[noitemsep]
    \item Left-right consistency check enabled: discards pixels where
          the left-to-right and right-to-left disparity estimates disagree
          beyond a threshold.
    \item Subpixel mode enabled: disparity encoded as 16-bit values with
          5 fractional bits (divide by 32 on host to recover pixel disparity).
    \item 7$\times$7 median filter for speckle suppression.
    \item Runtime-configurable confidence threshold via an \texttt{XLinkIn}
          stream, adjustable without restarting the pipeline.
\end{itemize}

\subsection{Queue Discipline}

Three output queues carry data from device to host.
The disparity queue is blocking (size~1): every new disparity frame must
be consumed before the next is accepted, driving the render rate.
Colour and rectified mono queues use \texttt{tryGet} (non-blocking) so
the host always has the latest frame without stalling.

\subsection{Disparity Pre-Processing}

Before uploading disparity to the GPU, the host applies three steps:

\begin{enumerate}[noitemsep]
    \item \textbf{Speckle filter:} \texttt{cv::filterSpeckles} with max
          blob size 150~px and max disparity jump 16.
          Removes isolated incorrect matches from rain or specular surfaces.
    \item \textbf{Subpixel conversion:} \texttt{convertTo(CV\_32F, 1.0/32.0)}
          recovers true pixel disparity from the RAW16 subpixel encoding.
    \item \textbf{IIR temporal filter:}
          $d_{\mathrm{out}} = 0.3 \cdot d_{\mathrm{new}} + 0.7 \cdot d_{\mathrm{prev}}$,
          with $d_{\mathrm{prev}}$ preserved unchanged where
          $d_{\mathrm{new}} = 0$.
          The zero-value guard is critical: without it, zero disparity
          from invalid pixels bleeds into adjacent valid estimates over
          time, growing the hole on every frame.
\end{enumerate}

\begin{figure}[h]
\centering
\figplaceholder{Disparity map before and after pre-processing.
  Left: raw subpixel disparity (RAW16, scaled for display). Right:
  after speckle filter and IIR temporal smoothing. Nearer objects
  appear brighter; invalid stereo regions appear as black holes
  filled by the JFA pass.}
\caption{OAK-D Lite disparity pre-processing.}
\label{fig:disparity_map}
\end{figure}

\section{Right-Disparity Construction}

\texttt{VWINDOW\_RIGHT\_CS} requires a disparity value for right-camera
pixels so that it can classify which right-camera pixels correspond to
genuinely disoccluded regions.
This is constructed on the CPU per scanline before the GPU pass:

\begin{enumerate}[noitemsep]
    \item \textbf{Forward pass}: for each left pixel $(x_l, y)$ with
          disparity $d$, write $d$ at right pixel
          $x_r = \mathrm{round}(x_l - d)$, keeping the maximum when
          multiple left pixels map to the same right pixel
          (nearest surface wins).
    \item \textbf{Backward fill}: scan right-to-left; fill zero
          (disoccluded) pixels with the last nonzero value to the right.
          Background depth propagates leftward into holes.
\end{enumerate}

% =====================================================================
%  CHAPTER 6 — THE DIBR PIPELINE: HEAD TRACKING
% =====================================================================
\chapter{The DIBR Pipeline: Head Tracking}

\section{Architecture Evolution}

The face tracker went through four complete architectures before reaching
the final design.
Table~\ref{tab:tracker_history} summarises the sequence and failure mode
of each.

\begin{table}[h]
\centering
\caption{Face tracker architecture evolution}
\label{tab:tracker_history}
\begin{tabular}{@{}p{4cm}p{7.5cm}@{}}
\toprule
Architecture & Failure mode / reason for replacement \\
\midrule
Lucas-Kanade optical flow
  & Drifted without periodic re-detection; anchor
    lost in low-texture frames. \\
OpenCV DNN (face + landmark)
  & Inference too slow on the Cortex-A57 CPU;
    consumed most of the available compute budget. \\
MOSSE + Haar Cascade
  & Achieved low CPU load, but provides only 2D
    screen coordinates.
    Metric depth is required to correctly scale
    parallax in the virtual window model. \\
\textbf{Haar + TRT FaceMesh}
  & \textbf{Final architecture.}
    Haar cascade for bootstrapping; TRT FaceMesh
    on iris-centred crop for subsequent frames;
    iris diameter provides metric depth. \\
\bottomrule
\end{tabular}
\end{table}

\subsection{MOSSE + Haar Cascade}

This was the first architecture to achieve acceptable CPU load on the
Jetson Nano.
A Haar cascade face detector ran every $N$ frames to anchor the tracker;
a Minimum Output Sum of Squared Error (MOSSE) correlation
filter~\cite{bolme2010mosse} maintained the bounding box between
detections.
A heavy IIR low-pass filter on the bounding box coordinates eliminated
jitter.
The face normal was computed from the bounding box aspect ratio and a
calibrated homography.

This approach was replaced because it cannot provide metric depth ---
the MOSSE bounding box gives only 2D screen position.
The virtual window model requires all three spatial coordinates of the
driver's head to correctly scale the parallax effect.

\section{Final Architecture: Haar + TRT FaceMesh}

\subsection{Detection and Tracking}

The face tracker runs on a dedicated thread consuming frames from the
viewer-facing USB webcam via a GStreamer pipeline
(\texttt{v4l2src} $\to$ \texttt{jpegdec} $\to$ \texttt{BGR} $\to$ \texttt{appsink},
\texttt{drop=true}, \texttt{max-buffers=1}).

The pipeline is:

\begin{enumerate}[noitemsep]
    \item \textbf{Haar cascade}~\cite{viola2004robust} on a 50\%-scale
          greyscale, histogram-equalised frame: scale factor 1.1, minimum
          4 neighbours, minimum face size 60$\times$60.
          The largest detection is selected.
    \item \textbf{Bounding box expansion and squaring}: 40\% margin added
          on each side, cropped to a square, resized to 192$\times$192
          RGB float [0,~1].
    \item \textbf{Bootstrap frame}: full-frame FaceMesh inference (first
          frame with a new detection) provides 478 landmarks.
          Subsequent frames crop around the iris centre for speed.
    \item \textbf{TensorRT FaceMesh engine}~\cite{tensorrt,kartynnik2019facemesh}
          (ONNX, FP16, 256~MB workspace): built from ONNX on first run,
          serialised to a \texttt{.trt} cache file, reloaded on subsequent
          runs.
    \item \textbf{Iris landmarks} (tensor output indices 468--477) provide
          the iris bounding circle for depth estimation via
          Equation~\ref{eq:iris_depth}.
    \item \textbf{IIR smoothing} on $(X, Y, Z)$ to reduce jitter.
\end{enumerate}

\begin{figure}[h]
\centering
\figplaceholder{Face tracker output. The 478-landmark FaceMesh is
  overlaid on the face crop; the iris bounding circle (landmarks
  468--477) is highlighted. The estimated 3D head position
  $(X, Y, Z)$ is shown in the overlay.}
\caption{Haar + TRT FaceMesh face tracker.}
\label{fig:face_tracker}
\end{figure}

\subsection{Coordinate Mapping}
\label{sec:coord_mapping}

The face tracker camera faces the viewer ($+Z$ toward camera), while
the OAK-D faces the scene ($+Z$ away from camera).
The tracker's $+X$ is the viewer's right, which is the scene's left
(OAK-D $-X$).
Two sign inversions are applied:

\begin{equation}
    h_x^{\mathrm{OAK}} = -h_x^{\mathrm{tracker}}, \quad
    h_z^{\mathrm{OAK}} = -h_z^{\mathrm{tracker}}
    \label{eq:coord_flip}
\end{equation}

Without these, the parallax direction is inverted --- moving the head
right makes the scene shift right instead of left.

\subsection{Calibration}

The tracker supports interactive recalibration (key \texttt{C}).
The current raw $(X, Y)$ position is stored as a reference offset;
subsequent output is the delta from this offset.
The driver's neutral forward-facing position therefore maps to $(0, 0)$
in the virtual window frame regardless of webcam mounting position.

% =====================================================================
%  CHAPTER 7 — THE DIBR PIPELINE: COMPUTE SHADERS
% =====================================================================
\chapter{The DIBR Pipeline: Compute Shaders}

\section{UNPROJECT\_CS}

This pass converts the uploaded disparity texture to a world-space
XYZ texture (RGBA32F).
For each pixel $(u, v)$ with disparity $d$, Equations~\ref{eq:depth_from_disparity}
and~\ref{eq:unproject} are applied to recover the world-space
coordinates.
Pixels with $d < 0.5$ (invalid disparity) write $w = 0$ and are skipped
in all subsequent passes.

\section{CLEAR\_CS}

The depth SSBO (one unsigned integer per output pixel) must be zeroed
each frame.
A dedicated compute shader with \texttt{local\_size\_x = 64} performs
this clear on the GPU, avoiding the overhead of uploading a
$\sim$1.2~MB zero buffer from the host each frame.

\section{VWINDOW\_DEPTH\_CS and VWINDOW\_COLOR\_CS}

For each valid world point $W$, a ray is cast from head position $H$ and
intersected with the display plane using Equation~\ref{eq:ray_plane}.
The intersection $Q$ is projected into display UV coordinates using the
display's right and up vectors.

The DEPTH pass resolves the many-to-one mapping when multiple world
points project to the same output pixel using an atomic depth test:

\begin{equation}
    \texttt{atomicMax}(\texttt{depth}[\text{dst}],\;
    \texttt{floatBitsToUint}(1/Z_W))
    \label{eq:atomic_depth}
\end{equation}

Storing $1/Z$ as float bits exploits a property of IEEE~754: positive
floats have the same bit-level ordering as unsigned integers, so
\texttt{atomicMax} on the bit pattern of $1/Z$ correctly selects the
nearest surface.
This avoids the need for floating-point atomic operations, which are
unavailable in OpenGL ES~3.1.

The COLOR pass reads each output pixel's depth value, back-projects it
to find the corresponding source pixel in the left camera image, and
writes the bilinear-sampled colour to the output texture.

\section{VWINDOW\_RIGHT\_CS}

When the virtual viewpoint shifts, regions occluded in the left camera
become visible.
The right rectified image provides colour data for these disocclusion
regions.
For each right-camera pixel with disparity $d_R$, the disocclusion filter
checks:

\begin{equation}
    d_R < d_L \cdot 0.7
    \label{eq:disocclusion_filter}
\end{equation}

where $d_L$ is the co-located left-camera disparity.
If satisfied, the pixel is a true disocclusion.
The right pixel is only written where \texttt{depth[dst] == 0}:
pixels already filled by the left-camera pass are not overwritten.

\section{Jump Flood Algorithm Hole Fill}

Remaining holes are filled using the Jump Flood
Algorithm~\cite{rong2006jfa}, which computes an approximate
nearest-seed Voronoi diagram in $\mathcal{O}(\log N)$ passes by
propagating seed coordinates outward at a step distance that halves
each iteration.

An initialisation pass (\texttt{JFA\_INIT\_CS}) marks each filled pixel
as a seed, storing its coordinate packed as \texttt{x | (y << 16)} in an
R32I image.
Each JFA pass checks the 9 neighbours at the current step distance and
adopts the nearest seed.
After $\lceil \log_2(\max(W, H)) \rceil$ passes, a gather pass copies
the nearest seed's colour into each hole pixel.

This replaces the earlier horizontal scan hole-fill, which only
propagated colour rightward, producing directional streaks at depth
discontinuities.

\begin{figure}[h]
\centering
\figplaceholder{JFA hole fill comparison. Left: output after
  left-camera colour splat and right-camera disocclusion fill --- unfilled
  holes appear black. Right: after JFA gather pass --- holes filled
  with the nearest valid seed colour, eliminating the directional streaking
  of the prior horizontal scan.}
\caption{Jump Flood Algorithm hole fill.}
\label{fig:jfa_holefill}
\end{figure}

\section{Display Pass}

The final fragment shader samples the hole-filled texture via a
fullscreen quad.
A \texttt{u\_scale} uniform applies NDC scaling to letterbox or
pillarbox the output to the correct aspect ratio.

% =====================================================================
%  CHAPTER 8 — KEY ENGINEERING CHALLENGES
% =====================================================================
\chapter{Key Engineering Challenges}

\section{Build System and DepthAI Dependencies}

Building the depthai-core library on L4T~R32 required the following
non-obvious steps:

\begin{itemize}[noitemsep]
    \item A Hunter package manager gate must precede the
          \texttt{project()} call --- non-standard for CMake.
    \item CMake minimum version required: 3.20, above the default L4T
          apt package.
    \item \texttt{PresetType::HIGH\_DENSITY} was removed from the API
          in the version available; replaced with a string-based preset.
    \item Frame data initially retrieved via \texttt{getCvFrame()}, which
          requires the \texttt{DEPTHAI\_OPENCV\_SUPPORT} compile flag;
          replaced with \texttt{getData()} and manual \texttt{cv::Mat}
          construction from the raw byte pointer.
\end{itemize}

\section{Camera Rig Orientation}

The OAK-D Lite was mounted inverted (rotated 180$^\circ$) for
installation reasons, inverting both the left and right mono streams and
the colour stream.
Compensation was applied in two places:

\begin{enumerate}[noitemsep]
    \item The GStreamer capture pipeline for the USB webcam head tracker
          was given a \texttt{videoflip flip-method=rotate-180} element.
    \item OAK-D frames were flipped in the host pre-processing step
          using \texttt{cv::rotate(frame, frame, cv::ROTATE\_180)}.
\end{enumerate}

\section{Coordinate System Consistency}

Three separate sign errors were introduced and corrected across
development.

\textbf{Parallax inversion.}
The face tracker camera faces the viewer, so its $+X$ is the viewer's
right --- opposite to the OAK-D scene frame.
The missing negation of $h_x$ (Equation~\ref{eq:coord_flip}) caused
parallax to invert: moving the head left made the scene shift left
instead of right.

\textbf{Viewer depth sign.}
The OAK-D faces into the scene along $+Z$, placing the viewer at
negative $Z$.
The tracker returns a positive viewer distance, requiring negation
before use as $h_z$.

\textbf{Display plane placement.}
The display was initially placed at $Z = 0$ (the OAK-D optical origin).
Scene points at 2--5~m projected onto this plane from a viewer at
$Z = -1$~m, producing a very narrow field of view.
Moving the display to $Z = -0.4$~m corrected the scale.
The final model places the display at $Z = 0$ with the OAK-D slightly
behind it, matching the physical geometry.

\section{Disocclusion Fill Correctness}

Achieving correct disocclusion fill required three iterations.

\textbf{Attempt 1: Pixel splatting.}
A 2$\times$2 splat on forward-warped pixels was intended to close
sub-pixel cracks.
It inadvertently filled holes with adjacent incorrect colours, producing
coloured fringing at depth edges.

\textbf{Attempt 2: Unconstrained right-camera write.}
\texttt{VWINDOW\_RIGHT\_CS} wrote right-camera colour into any output
pixel not yet filled.
This caused ghosting: right-camera pixels landed on already-correct
left-camera pixels due to the small baseline offset between the views.
Fix: only write where \texttt{depth[dst] == 0}.

\textbf{Attempt 3: Missing disocclusion filter.}
Ghosting persisted at depth discontinuities.
The right camera sees the same surface at a slightly different disparity,
so right pixels near edges were incorrectly treated as new content.
The disocclusion filter (Equation~\ref{eq:disocclusion_filter}) resolved
this: if both cameras see approximately the same depth, the right pixel
is from the same surface and is discarded.

\section{StereoDepthConfig State Reset}

Runtime stereo parameter updates are sent via an \texttt{XLinkIn} stream
as a \texttt{StereoDepthConfig} object.
Constructing a fresh config object to update one parameter silently resets
all other parameters to defaults.
Subpixel mode and left-right consistency check were both reset to off on
every keypress, causing disparity quality to degrade progressively after
any runtime adjustment.
The fix is to re-assert all desired settings in every configuration update:

\begin{lstlisting}[language=C++,
  caption={StereoDepthConfig update with explicit re-assertion.
    All flags must be set on every update; constructing a fresh config
    object silently resets omitted flags to their defaults.}]
StereoDepthConfig cfg;
cfg.setMedianFilter(enable_median ? KERNEL_7x7 : MEDIAN_OFF);
cfg.setLeftRightCheck(true);     // must re-assert every update
cfg.setSubpixel(true);           // must re-assert every update
configIn->send(cfg);
\end{lstlisting}

\section{IIR Temporal Filter Invalid-Sample Decay}

The IIR filter applied to the disparity map:

\begin{equation}
    d_{\mathrm{out}} = \alpha \cdot d_{\mathrm{new}} +
                       (1 - \alpha) \cdot d_{\mathrm{prev}}
    \label{eq:iir}
\end{equation}

was initially applied unconditionally.
Where the OAK-D returns zero disparity (no valid stereo match), zero was
blended in as if it were a real depth reading.
Over successive frames, valid depth estimates near persistent holes
decayed toward zero, growing the hole on every frame.
Fix: detect where $d_{\mathrm{new}} = 0$ and carry $d_{\mathrm{prev}}$
forward unchanged, preserving the last valid estimate.

% =====================================================================
%  CHAPTER 9 — RESULTS AND ANALYSIS
% =====================================================================
\chapter{Results and Analysis}

\section{Frame Rate}

Table~\ref{tab:fps_comparison} summarises pipeline frame rate at each
development phase.

\begin{table}[h]
\centering
\caption{Pipeline frame rate by development phase at 640$\times$480}
\label{tab:fps_comparison}
\begin{tabular}{@{}p{4cm}p{4.5cm}c@{}}
\toprule
Phase & Stereo method & FPS \\
\midrule
Phase 1: CSI + CPU SGBM    & CPU, 20--30~ms/frame    & $\approx$5 \\
Phase 1: CSI + CUDA BM     & Maxwell GPU, 10--15~ms  & $\approx$15 \\
Phase 3: OAK-D Lite        & Myriad X hardware, 0~ms host & 25--30 \\
\bottomrule
\end{tabular}
\end{table}

The dominant gain comes from offloading stereo to the Myriad X VPU,
freeing the Maxwell GPU and Cortex-A57 CPU entirely for the render
pipeline and head tracker.
At 0.5$\times$ processing scale ($320\times240$), frame rate exceeds
60~FPS with negligible visual quality loss for the A-pillar application,
which operates at coarser spatial frequencies than full resolution.

\section{Pipeline Latency Breakdown}

Table~\ref{tab:latency} shows the estimated per-stage latency for the
final system.
Because the OAK-D stereo pipeline and head tracker run on independent
threads, their latencies are pipelined rather than added sequentially.
The dominant end-to-end latency is the OAK-D stereo frame interval.

\begin{table}[h]
\centering
\caption{Estimated per-stage latency, Phase 3 system}
\label{tab:latency}
\begin{tabular}{@{}lcc@{}}
\toprule
Stage & Thread & Estimated latency \\
\midrule
OAK-D SGBM (hardware)         & OAK-D VPU (pipelined) & $\sim$30~ms \\
Host pre-processing (IIR, speckle) & OAK receiver        & $<$2~ms \\
GPU pipeline (UNPROJECT through JFA) & Render            & $\sim$3--5~ms \\
Face tracker (Haar + TRT FaceMesh)  & Dedicated thread  & $\sim$8--12~ms \\
\midrule
Total (render-limited)        & ---                   & $\sim$35~ms ($\approx$28~FPS) \\
\bottomrule
\end{tabular}
\end{table}

\section{Depth Quality}

The OAK-D Lite's 5-bit subpixel disparity (1/32~px resolution) reduces
depth quantisation substantially relative to integer-pixel matching.
At 4~m depth ($f = 381$~px, $B = 0.075$~m, integer disparity
$d \approx 7.1$~px), the depth change per pixel step is:

\begin{equation}
    \Delta Z = \frac{Z^2}{f \cdot B} \approx 0.56~\mathrm{m}
\end{equation}

With 5-bit subpixel, the minimum disparity step is $1/32$~px, giving:

\begin{equation}
    \Delta Z_{\mathrm{sub}} = \frac{Z^2 \cdot \delta d}{f \cdot B} =
    \frac{16 \times (1/32)}{381 \times 0.075} \approx 17.5~\mathrm{mm}
\end{equation}

This is a 32$\times$ improvement in depth resolution at 4~m, consistent
with the 5-bit fractional encoding.
Hardware benchmarks~\cite{arxiv2501} report $<$2\% absolute depth error
at 4~m for the OAK-D hardware under controlled conditions, corresponding
to $\pm$80~mm at this range.

Global shutter on the OV7251 sensors eliminates rolling-shutter skew,
which was a visible artefact with the CSI camera rig during fast lateral
head movements: rows captured at different times produced a laterally
sheared disparity map that broke the parallax effect momentarily.

\section{Parallax Effect}

The virtual window effect is most apparent at $\pm$50--100~mm of lateral
head displacement from the calibrated reference position.
Using the parallax shift formula derived from Equation~\ref{eq:unproject}:

\begin{equation}
    \delta u = \frac{f_x \cdot \Delta x}{Z}
\end{equation}

At 3~m depth with a 50~mm lateral displacement, the expected pixel shift
is $\delta u = 381 \times 0.05 / 3 \approx 6.4$~px.
At 100~mm displacement this rises to $\approx$12.7~px --- clearly
perceptible.
Near objects at 1~m shift proportionally more, providing strong depth cues.
No scaling constants are required: the virtual window model produces this
behaviour automatically from the physical geometry.

Table~\ref{tab:parallax_shift} shows predicted and observed pixel shifts
at representative depths and head displacements.

\begin{table}[h]
\centering
\caption{Predicted parallax pixel shift vs.\ depth and lateral displacement
         ($f_x \approx 381$~px)}
\label{tab:parallax_shift}
\begin{tabular}{@{}lccc@{}}
\toprule
Depth & $\Delta x = 30$~mm & $\Delta x = 60$~mm & $\Delta x = 100$~mm \\
\midrule
1~m  & 11.4~px & 22.9~px & 38.1~px \\
3~m  & 3.8~px  & 7.6~px  & 12.7~px \\
5~m  & 2.3~px  & 4.6~px  & 7.6~px  \\
\bottomrule
\end{tabular}
\end{table}

Shifts above $\approx$3~px are reliably perceptible in a stationary
scene.
At 3~m depth, the effect is clearly visible from 60~mm displacement; at
1~m, it is visible from as little as 10~mm.

\section{Remaining Artefacts}

\textbf{Large disocclusions.}
When the head moves more than approximately 60~mm, background regions
become visible with no colour data in either camera.
JFA fills these with the nearest valid colour, which is plausible for
small holes but visibly incorrect at high-contrast boundaries.

\textbf{Textureless regions.}
Passive stereo fails on uniform surfaces (plain walls, overcast sky).
The OAK-D confidence map is used to mask low-confidence disparity before
upload; textureless areas remain as holes filled by JFA.

\textbf{Iris depth variation.}
The fixed iris diameter constant introduces approximately $\pm$10\% depth
error across the adult population, manifesting as slight over- or
under-scaling of parallax magnitude.

\section{Display Calibration}

Direct calibration between the external OAK-D and the interior display is
geometrically difficult: the two devices face opposite directions and are
separated by the car body.
A \texttt{solvePnP}-based approach chaining transforms through an
intermediate reference marker was investigated but found to accumulate
excessive error (1--2$^\circ$ angular error produces centimetre-scale
positional offset at display distance).
The adopted procedure initialises the display-to-camera transform from
physical measurements and refines it by observing whether parallel scene
lines remain parallel in the output; the dominant uncertainty sources
(iris depth $\pm5$--10\%, tracker noise $\pm5$~mm) exceed the calibration
error at this stage, making a more rigorous procedure unnecessary.

% =====================================================================
%  CHAPTER 10 — CONCLUSIONS AND REFLECTIONS
% =====================================================================
\chapter{Conclusions and Reflections}

\section{Summary}

ClearVision demonstrates that a real-time, geometrically correct
A-pillar transparency system is achievable on embedded GPU hardware at
a cost accessible to a student project.
The system achieves 25--30~FPS at $640\times480$ on a Jetson Nano,
compared to approximately 5~FPS in the initial prototype, by offloading
stereo depth to the OAK-D Lite's Myriad X VPU and implementing the
render pipeline as a multi-pass OpenGL ES~3.1 compute shader chain.
The virtual window projection model produces physically correct parallax
with no empirical tuning constants.

Two Python proof-of-concept demonstrations --- head pose vector
visualisation and head-coupled off-axis projection --- validated the core
geometry and established design decisions before the full pipeline was
built, substantially reducing implementation risk.

\section{Key Technical Contributions}

\begin{itemize}[noitemsep]
    \item \textbf{Virtual window projection model}: replaces the
          ad-hoc disparity shift approach with geometrically correct
          ray-plane intersection.
          Parallax scale is physically derived, with no tuning constants.
    \item \textbf{JFA hole fill}: $\mathcal{O}(\log N)$ 2D nearest-seed
          Voronoi fill replaces $\mathcal{O}(N)$ horizontal scan,
          eliminating directional streaking.
    \item \textbf{\texttt{floatBitsToUint} depth ordering}: enables
          \texttt{atomicMax} depth testing in OpenGL ES~3.1 without
          floating-point atomic support, exploiting the IEEE~754 float
          bit ordering.
    \item \textbf{IIR invalid-sample guard}: prevents valid depth from
          decaying into adjacent zero-disparity holes over time.
    \item \textbf{Iris-based metric depth}: uses the known iris diameter
          (11.7~mm) as a depth cue, avoiding a depth sensor on the
          driver-facing side of the system.
\end{itemize}

\section{Limitations}

\textbf{Large disocclusions.}
JFA nearest-neighbour fill is plausible only for small holes
(roughly 30~px).
For larger disocclusions, background inpainting or temporal reprojection
would be required.
This is a fundamental constraint of single-viewpoint DIBR.

\textbf{Textureless regions.}
Passive stereo fails on uniform surfaces.
Active stereo (structured light) or a time-of-flight sensor would
provide reliable depth in such scenes at substantially higher hardware
cost.

\textbf{Single viewer.}
The system is calibrated to one driver's head position.
Multi-viewer support would require time-multiplexed rendering or a
lenticular display.

\textbf{Display calibration.}
The manual calibration procedure is sufficient for development but would
require a repeatable mechanical fixture for production deployment.

\section{Future Work}

\begin{itemize}[noitemsep]
    \item \textbf{Temporal reprojection}: warp the previous frame's
          output into the current frame using the inter-frame head
          displacement, seeding holes before JFA.
          This would reduce visible artefacts for large disocclusions.
    \item \textbf{Joint bilateral depth upsampling}: upsample the 480P
          depth map to match the IMX214 colour resolution (up to 1080P),
          guided by colour edges, before the UNPROJECT pass.
    \item \textbf{IMU-assisted display pose}: use the OAK-D Lite's
          onboard Bosch BMI270 and a display-mounted IMU to automate
          the rotation component of the display-to-camera calibration.
\end{itemize}

\section{Reflections}

The project demonstrated that hardware-accelerated stereo depth is now
accessible at consumer price points.
The OAK-D Lite substantially lowered the barrier to entry: hardware
SGBM, subpixel disparity, global shutter, EEPROM-stored intrinsics, and
a single USB-C cable replaced what previously required a CUDA-capable GPU,
a manual calibration rig, and careful CSI ribbon routing.

The majority of development time was spent not on the core DIBR rendering
but on peripheral problems: build system compatibility, tracker reliability,
and coordinate system consistency.
These are typical of projects that integrate multiple independently
developed components (DepthAI, TensorRT, OpenGL ES, Qt) on a constrained
embedded platform.
Documenting the specific failure modes encountered provides a concrete
reference for similar work on the same hardware stack.

The virtual window model, once correctly implemented, required no further
calibration.
Physically correct parallax --- including correct vertical parallax,
depth-dependent magnitude, and natural near-object motion --- made the
A-pillar elimination effect substantially more convincing than the earlier
disparity-shift model.

% =====================================================================
%  BIBLIOGRAPHY
% =====================================================================
\bibliographystyle{ieeetr}
\bibliography{references}

% =====================================================================
%  APPENDICES
% =====================================================================
\appendix

\chapter{Stereo Baseline Selection}
\label{app:baseline}

The stereo baseline determines both depth accuracy and the maximum
synthesisable viewpoint shift.
For the A-pillar application, both constraints converge on the same range.

The A-pillar obstruction subtends approximately 80--120~mm of lateral
space at the driver's eye position (measured on a typical passenger
vehicle, driver seated 600--700~mm from the pillar).
The synthesised viewpoint shift must span this width.
A narrower baseline would not produce sufficient shift; a wider baseline
would shift near-pillar objects beyond the obstruction width, breaking
the window illusion.

The OAK-D Lite's 75~mm fixed baseline lies comfortably within this
range.

A further constraint: extending the baseline too far increases
disocclusion region size, worsening visible artefacts from the JFA fill.
For DIBR specifically, the recommended stereo baseline is approximately
5--10\% of the target scene depth --- for 3--5~m scenes, this gives
150--500~mm.
The A-pillar geometry constraint at 80--120~mm is therefore the binding
constraint, not depth accuracy.

\chapter{Key Source File Reference}
\label{app:source}

\begin{table}[h]
\centering
\caption{Principal source files}
\begin{tabular}{@{}p{5.5cm}p{7cm}@{}}
\toprule
File & Role \\
\midrule
\texttt{src/view\_synthesis.cpp}
  & Qt OpenGL widget; compute shader dispatch; main render loop \\
\texttt{src/oak\_receiver.cpp}
  & DepthAI pipeline; OAK-D polling thread \\
\texttt{src/face\_tracker.cpp}
  & Haar + TRT FaceMesh pipeline; iris depth estimation \\
\texttt{src/face\_tracker.h}
  & Tuning constants (\texttt{FT\_SMOOTH}, iris diameter) \\
\texttt{src/calibration.cpp}
  & Legacy manual stereo calibration --- retained for reference \\
\texttt{demos/head\_coupled\_offaxis.py}
  & Off-axis projection proof-of-concept \\
\texttt{demos/face\_landmarker\_demo.py}
  & Head pose vector proof-of-concept \\
\bottomrule
\end{tabular}
\end{table}

\end{document}
